\section{Quantitative empirical balance of the CUE labels}
\label{sec:analytic-cue-empirical-balance}

NN-A003 concerns a uniformly marked eigenangle. The following
statement controls the actual fraction of labels in a single
sample, while retaining the same correlated exterior phase.

\begin{theorem}[NN-A031: empirical half proportions in mean square]
\label{thm:analytic-cue-empirical-balance}
Fix $r_0<1$. For each $N$, use the CUE labels of NN-A003 with
arbitrary $0\le r_N\le r_0$ and arbitrary offset $a_N\in\mathbb R$:
\[
 W_{N,i}=i-\frac{N\theta_i}{2\pi}-a_N
          -\frac1\pi\operatorname{Im}\sum_j
                   \log(1-r_Ne^{i(\theta_i-\theta_j)}),
 \qquad \eta_{N,i}=\operatorname{sgn}\sin(\pi W_{N,i}).
\]
Almost surely none of these labels is zero. There is a constant
$C_{r_0}$, independent of $N,r_N,a_N$, such that for $N\ge3$,
\begin{equation}\label{eq:analytic-cue-empirical-rate}
 \mathbb E\left[\left(\frac1N\sum_{i=1}^N\eta_{N,i}\right)^2\right]
                          \le\frac{C_{r_0}}{\log N}.
\end{equation}
In particular, the observed positive-label fraction $L_N$ satisfies
\[
 \mathbb E[(L_N-1/2)^2]\le\frac{C_{r_0}}{4\log N},\qquad
 \mathbb P(|L_N-1/2|>\varepsilon)
                   \le\frac{C_{r_0}}{4\varepsilon^2\log N}.
\]
Thus $L_N\to1/2$ in $L^2$ and in probability. No optimality of
the logarithmic error bound, almost-sure convergence along all
$N$, or arithmetic transfer is asserted.
\end{theorem}

\begin{proof}
Write $P=P_N$ for the CUE projection in
$L^2(S^1,d\theta/(2\pi))$, and let $I,J$ be two distinct ordered
uniform indices sampled after the matrix. All constants below
are independent of the offset and of $r\in[0,r_0]$.

\emph{An exact uniform marginal modulo one.}
Under simultaneous rotation by $\alpha$, the exterior phase at
a carried marked point is unchanged. Its coordinate changes by
$\alpha$ modulo $2\pi$, and its new index differs from its old
index by an integer. Consequently its $W$ coordinate modulo one
changes by $-N\alpha/(2\pi)$. The marked Haar ensemble is
rotation invariant, so for every nonzero integer $h$,
\[
 \mathbb E e^{2\pi ihW_{N,I}}
 =e^{-ihN\alpha}\mathbb E e^{2\pi ihW_{N,I}}
 \quad\hbox{for every }\alpha.
\]
These expectations vanish. Hence $W_{N,I}\bmod1$ is exactly
uniform for every $N$. In particular a uniformly marked label
is almost surely nonzero, and a finite union gives the same
assertion for all labels in a sample. For any integrable
one-periodic function $g$ we have the exact identity
\begin{equation}\label{eq:analytic-cue-exact-modone}
             \mathbb E g(W_{N,I})=\int_0^1g(u)\,du.
\end{equation}

\emph{Two-Palm odd Fourier modes.}
Condition on the marked eigenangles $x\ne y$. The remaining
points form a projection determinantal process with projection
$Q_{N,x,y}$ obtained by removing the two-dimensional span of
the evaluation vectors at $x,y$ from $\operatorname{ran}P$.
This follows by the Schur complement of their two rows and
columns in the CUE density; the two evaluation vectors are
linearly independent for $N\ge2$ and $x\ne y$.

Let $k,l$ be odd integers. Their index parities combine to the
parity of the number of remaining eigenangles in the arc
between $x$ and $y$ that does not cross the coordinate origin.
For example, when $x<y$, the two counts in $(0,x)$ and $(0,y)$
have sum congruent modulo two to the count in $(x,y)$; the
contribution from the two marked points is a constant sign.
Define
\[
 j_{x,y}(\theta)=(-1)^{\mathbf1_{(\min(x,y),\max(x,y))}(\theta)},
 \qquad g_x(\theta)=\operatorname{Im}\log(1-re^{i(x-\theta)}),
\]
and
\[
 v(\theta)=e^{-ikg_x(\theta)-ilg_y(\theta)},\qquad u=j_{x,y}v.
\]
The self- and cross-contributions of the marked points to the
exterior phases are fixed scalars of modulus one. The determinant
identity in NN-A003 therefore gives
\begin{equation}\label{eq:analytic-cue-two-palm-mode}
 \left|\mathbb E\left[e^{i\pi(kW_{N,I}+lW_{N,J})}
                     \mid x,y\right]\right|
   =\left|\det(Q_{N,x,y}uQ_{N,x,y})\right|.
\end{equation}
The determinant is on the range of $Q_{N,x,y}$.

For a unimodular multiplier $w$, retain the notation
\[
 F_N(w)=\|(1-P)wP\|_{\rm HS}^2
       =\sum_{n\in\mathbb Z}\min(N,|n|)|\widehat w(n)|^2.
\]
Expanding the two off-diagonal blocks of the commutator gives
$\|[P,w]\|_{\rm HS}^2=2F_N(w)$: both blocks have squared norm
$N-\|PwP\|_{\rm HS}^2$. Since
$[P,jv]=[P,j]v+j[P,v]$, the triangle inequality gives
\begin{equation}\label{eq:analytic-cue-smooth-multiplier}
       \sqrt{F_N(jv)}\ge\sqrt{F_N(j)}-\sqrt{F_N(v)}.
\end{equation}
The smooth multiplier has derivative bounded by
\[
 \|v'\|_2\le C_0(|k|+|l|),\qquad C_0=\frac{r_0}{1-r_0}.
\]
By Cauchy--Schwarz and Parseval,
\begin{equation}\label{eq:analytic-cue-smooth-leakage}
 F_N(v)\le\sum_n|n||\widehat v(n)|^2
     \le\|v'\|_2\|v\|_2
     \le C_0(|k|+|l|).
\end{equation}

Let $d(x,y)\in(0,\pi]$ be their circular separation and let
$\ell=|x-y|$. The arc multiplier has the exact coefficients
\[
 |\widehat j_{x,y}(n)|^2
       =\frac{2-2\cos(n\ell)}{\pi^2n^2}\quad(n\ne0).
\]
Consequently
\[
 F_N(j_{x,y})\ge\frac4{\pi^2}
       \left(H_N-\sum_{n=1}^N\frac{\cos(n\ell)}n\right).
\]
Summation by parts bounds the tail of
$\sum_{n\ge1}e^{in\ell}/n$ by
$2/(N|1-e^{i\ell}|)$. Its real infinite sum is
$-\log|1-e^{i\ell}|$. Since
$|1-e^{i\ell}|\ge2d(x,y)/\pi$, for
$d(x,y)\ge N^{-1/2}$ we obtain, uniformly,
\[
 F_N(j_{x,y})\ge\frac2{\pi^2}\log N-O(1)
                         \ge\frac1{\pi^2}\log N
\]
for all sufficiently large $N$.

Put
\[
 c_* =\frac1{64\pi^2(1+C_0)},\qquad
 M=\lfloor c_*\log N\rfloor,
\]
and assume $N$ is sufficiently large that $M\ge2$. If $k,l$ are
odd and $|k|,|l|\le2M$, then
\[
 F_N(v)\le4C_0M\le\frac{\log N}{16\pi^2}
                         \le\frac1{16}F_N(j_{x,y}).
\]
Equation \eqref{eq:analytic-cue-smooth-multiplier} yields
$F_N(u)\ge(9/16)F_N(j_{x,y})\ge(\log N)/(2\pi^2)$.
Removing a rank-one projection changes the leakage by at most
one, as proved in NN-A003; hence the two-Palm leakage is at
least $F_N(u)-2$. The determinant contraction from NN-A003 gives
\[
 |\det(Q_{N,x,y}uQ_{N,x,y})|
                 \le eN^{-\kappa},\qquad
                      \kappa=\frac1{4\pi^2},
\]
whenever $d(x,y)\ge N^{-1/2}$.

The joint density of the ordered marked pair, relative to
$dx\,dy/(2\pi)^2$, is
\[
 \frac{N^2-|P_N(x,y)|^2}{N(N-1)}\le\frac{N}{N-1}\le2.
\]
The excluded region $d(x,y)<N^{-1/2}$ has measure
$O(N^{-1/2})$, and a conditional characteristic function has
absolute value at most one there. Integration therefore proves
the uniform estimate
\begin{equation}\label{eq:analytic-cue-uniform-odd-pair}
 \left|\mathbb E e^{i\pi(kW_{N,I}+lW_{N,J})}\right|
                  \le C N^{-\kappa}
 \quad\left(\begin{array}{c}k,l\text{ odd},\\|k|,|l|\le2M\end{array}\right).
\end{equation}
No independence of the exterior phases has been used.

\emph{A positive finite Fourier approximation to the labels.}
On $\mathbb R/2\mathbb Z$, with normalized measure $du/2$, let
$f(u)=\operatorname{sgn}\sin(\pi u)$ and normalize the positive
Jackson kernel
\[
 J_M(u)=\frac1{Z_M}
       \left|\frac{\sin(M\pi u/2)}{\sin(\pi u/2)}\right|^4,
 \qquad \int_{-1}^1J_M(u)\frac{du}{2}=1.
\]
The quotient has its continuous value $M$ at zero. Elementary
sine bounds give $Z_M\asymp M^3$ and
\[
 J_M(u)\le CM\min\{1,(M|u|)^{-4}\},\qquad
 \int_{-1}^1|u|J_M(u)\frac{du}{2}\le\frac CM.
\]
For the lower bound on $Z_M$, restrict to $|u|\le1/M$ and use
$|\sin(M\pi u/2)|\ge M|u|$ and
$|\sin(\pi u/2)|\le\pi|u|/2$. The upper and moment bounds follow
by splitting at $1/M$.

Set $q_M=f*J_M$. This is a real trigonometric polynomial of
degree at most $2(M-1)$, with $|q_M|\le1$. Because
$f(u+1)=-f(u)$, it has only odd Fourier modes. Moreover,
\begin{equation}\label{eq:analytic-cue-jackson-error}
 \|f-q_M\|_{L^1(\mathbb R/2\mathbb Z)}\le\frac CM,
 \qquad \sum_k|\widehat q_M(k)|\le C\log(M+1).
\end{equation}
The first estimate follows by averaging
$\|f(\cdot-u)-f\|_1\le2|u|$ against $J_M$. For the second use
$|\widehat J_M(k)|\le1$ and
$|\widehat f(k)|=2/(\pi|k|)$ for odd $k$.

The function $|f-q_M|$ is one-periodic. Thus
\eqref{eq:analytic-cue-exact-modone} makes its expected value at
the marked $W$ exactly its circle integral, without an
asymptotic equidistribution error. Since all labels and $q_M$
values have absolute value at most one,
\[
 \left|\mathbb E[\eta_{N,I}\eta_{N,J}]
       -\mathbb E[q_M(W_{N,I})q_M(W_{N,J})]\right|
                                             \le\frac CM.
\]
Using the odd-mode estimate
\eqref{eq:analytic-cue-uniform-odd-pair} in the finite Fourier
expansion gives
\[
 |\mathbb E[\eta_{N,I}\eta_{N,J}]|
       \le\frac CM+C N^{-\kappa}\log^2(M+1).
\]
Finally the diagonal terms of the empirical square contribute
$1/N$, so
\[
 \mathbb E\left[\left(\frac1N\sum_i\eta_{N,i}\right)^2\right]
 \le\frac1N+\frac CM+C N^{-\kappa}\log^2(M+1)
                              =O_{r_0}(1/\log N).
\]
Increasing the constant covers the finitely many smaller $N$.
The assertions about $L_N=(1+N^{-1}\sum_i\eta_{N,i})/2$ follow
from this identity and Markov's inequality.
\end{proof}

\paragraph{Geometric translation.}
Identify $\operatorname{ran}P_N$ with the restrictions to the unit
circle of $H^0(\mathbb P^1,\mathcal O(N-1))$, carrying its Haar
Hermitian norm. At two distinct marked points $z_x,z_y$, polynomial
division gives the exact evaluation sequence
\[
 0\longrightarrow (z-z_x)(z-z_y)\mathbb C[z]_{\le N-3}
 \longrightarrow\mathbb C[z]_{\le N-1}
 \longrightarrow\mathbb C^2\longrightarrow0.
\]
The two-Palm projection is the orthogonal projection onto this
vanishing subspace. Its exterior determinant controls the marked
Fourier statistics; the rank-two correction and the commutator
norm make that control quantitative. The representation-ring
exterior-power formalism is SGA~6, Expos\'e~V, Example~2.2(b) and
Lemma~2.2.1. These actual evaluation spaces and Haar projections
give a probabilistic balance law. They do not provide an
arithmetic comparison map or a Weil purity argument.
